% =====================================================================
% Mathematical dictionary
%
% The shared definitions used across every worksheet. This file is both a
% reference in its own right and the source the generator reads: when a
% worksheet's vocabulary key uses one of these words, the wording below is the
% one that appears on it, so that a pupil meets the same explanation every time.
%
% TO ADD OR CHANGE AN ENTRY
%
%   \dictentry{numerator}{the number above the line in a fraction}
%
% Write the definition as you would say it to a pupil who has not met the word:
% plain words, one sentence, no symbols, and never the word itself inside its
% own definition. Keep it short enough to sit in a table cell.
%
% Other forms of the word are found automatically - "numerators" and
% "Numerator" both match the entry above, and so do regular verb endings. Where
% a form is irregular enough that it would not be found, list it first:
%
%   \dictentry[vertices]{vertex}{a corner where two or more edges meet}
%
% Changing a wording here rebuilds the vocabulary keys that use that word, and
% only those, the next time the generator runs.
%
% Entries are grouped by topic for reading, but the order does not matter to
% the generator. Keep each group alphabetical.
% =====================================================================
\documentclass[10pt,twoside]{article}
\usepackage[a4paper,margin=16mm,columnsep=8mm]{geometry}
\usepackage{multicol}
\usepackage{xcolor}
\usepackage{fancyhdr}

\definecolor{headword}{RGB}{0,128,0}

% One entry. The optional argument lists extra forms of the word, and is shown
% after the headword so a reader can see what else counts as the same word.
\newcommand{\dictentry}[3][]{%
  \par\addvspace{2pt}%
  \noindent\textbf{\textcolor{headword}{#2}}%
  \if\relax\detokenize{#1}\relax\else{\small\textit{ (also #1)}}\fi
  \\ #3\par
}

\newcommand{\dicttopic}[1]{%
  \par\addvspace{6pt}%
  {\normalsize\bfseries #1}\par
  \vspace{1pt}\hrule\vspace{3pt}%
}

\pagestyle{fancy}
\fancyhf{}
\fancyhead[L]{\small Mathematical dictionary}
\fancyhead[R]{\small\thepage}
\renewcommand{\headrulewidth}{0.4pt}

\setlength{\parindent}{0pt}

\begin{document}

{\LARGE\bfseries Mathematical dictionary}\par
\vspace{2pt}
{\small The words a mathematics lesson depends on, explained in plain English.}\par
\vspace{6pt}\hrule\vspace{8pt}

\begin{multicols}{3}
\small

\dicttopic{Number}

\dictentry{decimal}{a number written with a point, where the digits after the point are parts of a whole}
\dictentry{denominator}{the number below the line in a fraction, showing how many equal parts the whole is split into}
\dictentry{difference}{the answer when one number is subtracted from another}
\dictentry{digit}{a single figure from 0 to 9 that a number is written with}
\dictentry{equivalent}{having the same value, even though it is written differently}
\dictentry{estimate}{a sensible answer worked out quickly from rounded numbers}
\dictentry{factor}{a whole number that divides exactly into another number}
\dictentry{fraction}{a number written as one whole split into equal parts, with a number of those parts taken}
\dictentry[HCF]{highest common factor}{the largest whole number that divides exactly into two or more numbers}
\dictentry{integer}{a whole number, which may be positive, negative or zero}
\dictentry[LCM]{lowest common multiple}{the smallest number that two or more numbers all divide into exactly}
\dictentry[multiplies, multiplied, multiplying, multiplication]{multiply}{to add a number to itself a given number of times}
\dictentry{multiple}{a number you get by multiplying another number by a whole number}
\dictentry{numerator}{the number above the line in a fraction, showing how many equal parts are taken}
\dictentry{percentage}{a number of parts out of a hundred}
\dictentry{place value}{the value a digit has because of where it sits in a number}
\dictentry[powers]{power}{how many times a number is multiplied by itself}
\dictentry{prime}{a whole number greater than 1 that divides only by itself and 1}
\dictentry{product}{the answer when two or more numbers are multiplied}
\dictentry{quotient}{the answer when one number is divided by another}
\dictentry[reciprocals]{reciprocal}{the number you multiply by to get 1, found by turning a fraction upside down}
\dictentry[recurs, recurring]{recur}{to repeat the same digits forever after the decimal point}
\dictentry[rounds, rounded, rounding]{round}{to replace a number with a nearby simpler one}
\dictentry[simplifies, simplified, simplifying]{simplify}{to write something in its smallest or plainest form without changing its value}
\dictentry[square root]{root}{a number that gives another number when multiplied by itself a given number of times}
\dictentry{sum}{the answer when numbers are added together}
\dictentry{surd}{a root that cannot be written exactly as a fraction, so it is left in root form}

\dicttopic{Ratio and proportion}

\dictentry{proportion}{a relationship where two amounts change by the same factor}
\dictentry{ratio}{a way of comparing two or more amounts, showing how much of one there is for each of the other}
\dictentry{scale factor}{the number every length is multiplied by when a shape is enlarged}
\dictentry{unitary}{a method that first finds the value of one part or one item}

\dicttopic{Algebra}

\dictentry{coefficient}{the number in front of a letter, showing how many of it there are}
\dictentry[expands, expanded, expanding]{expand}{to multiply out brackets so the expression has no brackets left}
\dictentry{expression}{a group of numbers and letters joined by operations, with no equals sign}
\dictentry[factorisation, factorize]{factorise}{to write an expression as a product, usually by putting brackets back in}
\dictentry{formula}{a rule written with letters, showing how to work one quantity out from others}
\dictentry{gradient}{how steep a line is, found by how far it rises for each step across}
\dictentry{identity}{a statement that is true for every value of the letters in it}
\dictentry[indices]{index}{the small raised number showing what power something is raised to}
\dictentry{inequality}{a statement that one quantity is larger or smaller than another}
\dictentry{intercept}{the point where a graph crosses an axis}
\dictentry{linear}{making a straight line when drawn as a graph}
\dictentry{quadratic}{an expression or equation whose highest power is a square}
\dictentry[rearranges, rearranged, rearranging]{rearrange}{to change the order of a formula so a different letter is on its own}
\dictentry{sequence}{a list of numbers that follows a rule}
\dictentry[substitutes, substituted, substituting, substitution]{substitute}{to put a number in place of a letter}
\dictentry{term}{one part of an expression or one number in a sequence}
\dictentry{nth term}{a rule that gives any number in a sequence from its position}
\dictentry{variable}{a letter standing for a quantity that can change}

\dicttopic{Geometry}

\dictentry{adjacent}{next to, and sharing a side or a corner}
\dictentry{arc}{part of the curved edge of a circle}
\dictentry{area}{the amount of space inside a flat shape}
\dictentry{bearing}{a direction given as an angle clockwise from north, written with three figures}
\dictentry{bisector}{a line that cuts an angle or another line exactly in half}
\dictentry{chord}{a straight line joining two points on the edge of a circle}
\dictentry{circumference}{the distance all the way round a circle}
\dictentry{congruent}{exactly the same shape and size}
\dictentry{diameter}{a straight line across a circle through its centre}
\dictentry[edges]{edge}{a straight line where two faces of a solid meet}
\dictentry[faces]{face}{a flat surface of a solid shape}
\dictentry{hypotenuse}{the longest side of a right-angled triangle, opposite the right angle}
\dictentry{parallel}{always the same distance apart and never meeting}
\dictentry{perimeter}{the total distance around the edge of a shape}
\dictentry{perpendicular}{meeting at a right angle}
\dictentry{polygon}{a flat shape with straight sides}
\dictentry{radius}{the distance from the centre of a circle to its edge}
\dictentry{sector}{a slice of a circle between two radii}
\dictentry{segment}{the part of a circle cut off by a chord}
\dictentry{similar}{the same shape but a different size}
\dictentry{tangent}{a line that touches a curve at exactly one point}
\dictentry{transformation}{a change to a shape's position, size or orientation}
\dictentry{vector}{a quantity with both size and direction}
\dictentry[vertices]{vertex}{a corner where two or more edges meet}
\dictentry{volume}{the amount of space inside a solid shape}

\dicttopic{Statistics and probability}

\dictentry{average}{a single value chosen to stand for a whole set of data}
\dictentry[axes]{axis}{one of the lines a graph is measured against}
\dictentry{correlation}{how closely two sets of data are related}
\dictentry{cumulative}{adding up as you go along, so each total includes the ones before}
\dictentry{data}{the set of values collected}
\dictentry{frequency}{how many times something happens}
\dictentry{mean}{the average found by adding all the values and dividing by how many there are}
\dictentry{median}{the middle value when the data is put in order}
\dictentry{mode}{the value that appears most often}
\dictentry{outlier}{a value far away from all the others}
\dictentry{probability}{how likely something is, given as a number between 0 and 1}
\dictentry{quartile}{a value marking one quarter of the way through ordered data}
\dictentry{range}{the difference between the largest and smallest values}
\dictentry{sample}{a smaller group chosen to stand for a larger one}

\dicttopic{Words not yet checked}

% ---- Words met on a worksheet, not yet checked by anybody ----
%
% The generator adds a word here the first time one of its vocabulary
% keys uses a word this file does not define, along with the definition
% that key was written with. Nobody has read these yet, so they are the
% entries most worth checking.
%
% Rewording one here changes it everywhere: every vocabulary key that
% uses that word is rebuilt to match, in English and in every language
% the repository is translated into. So this is the place to fix an
% explanation that reads badly, rather than the sheet it appeared on.
%
% Move an entry up into the topic it belongs to whenever you like. The
% generator does not mind where in the file an entry lives - it reads
% the word and the definition and nothing else.
% -------------------------------------------------------------
\dictentry{absolute value}{the distance of a number from zero, ignoring whether it is positive or negative.}
\dictentry{addition}{combining two or more numbers to find their total.}
\dictentry{altitude}{height above sea level or another fixed reference level.}
\dictentry{angle}{the amount of turn between two straight lines that meet at a point}
\dictentry{annular}{Describing a shape made by the space between two concentric circles.}
\dictentry{apex angle}{the angle at the point of a triangle opposite its base.}
\dictentry{arccos}{The operation that finds an angle when its cosine is known.}
\dictentry{arctan}{The operation that finds an angle when its tangent is known.}
\dictentry{area model}{a rectangle diagram split into smaller rectangles, used to show how to multiply fractions.}
\dictentry{arithmetic}{the branch of mathematics that deals with adding, subtracting, multiplying and dividing numbers.}
\dictentry{ascending order}{arranged from smallest to largest}
\dictentry{average speed}{total distance travelled divided by total time taken.}
\dictentry{bar model}{a diagram made of rectangles that shows the parts and total of an equation.}
\dictentry{base}{the side of a shape chosen as the bottom for working out its area.}
\dictentry{basis point}{a unit equal to one hundredth of one percent, used for small changes in rates.}
\dictentry{bimodal}{having two values that appear most often in a list of numbers.}
\dictentry{bodmas}{a rule that says which operation to do first in a calculation.}
\dictentry{bond}{a type of loan to a government or company that pays regular money and is paid back later.}
\dictentry{box method}{a way of multiplying or expanding by splitting a rectangle into smaller parts.}
\dictentry{bracket}{one of a pair of curved symbols, ( ), used to group numbers or letters.}
\dictentry{brackets}{A pair of signs used in mathematics to group together a part of an expression.}
\dictentry{cancel}{to remove a shared number from the top and bottom of a fraction before multiplying.}
\dictentry{cardinal directions}{The four main compass points: North, South, East and West.}
\dictentry{category}{a group of non-numerical items, for example a colour, used as data.}
\dictentry{centre}{the fixed middle point of a circle}
\dictentry{circle}{a round flat shape whose edge is always the same distance from its centre}
\dictentry{circular}{round in shape, with every point on its edge the same distance from the centre.}
\dictentry{column vector}{a vector written vertically with one number above another.}
\dictentry{common denominator}{a bottom number shared by two or more fractions, used for adding or comparing them.}
\dictentry{common difference}{the fixed number added to each term to get the next term}
\dictentry{common factor}{a number that divides exactly into two or more other numbers.}
\dictentry{common ratio}{the fixed number used to create each next term in a geometric sequence}
\dictentry{commutativity}{the rule that swapping two numbers before you add or times them leaves the answer the same}
\dictentry{component}{one of the numbers that make up a vector.}
\dictentry{compounding}{the process of earning profit on both the original amount and profit already earned.}
\dictentry{concentric}{Having the same centre.}
\dictentry{constant}{a fixed number that stands alone, like 3 in 3x+3.}
\dictentry{convergence}{the process of getting closer and closer to a particular value}
\dictentry{conversion}{a change of a value from one form, such as a fraction, to another, such as a percentage.}
\dictentry{convert}{to change a number from one form into another, such as a mixed number into an improper fraction}
\dictentry{approximately}{close to but not exactly; about.}
\dictentry{ascending}{arranged from smallest to largest.}
\dictentry{bus stop method}{a written method for division using a roof-like sign and a divisor.}
\dictentry{coordinate}{a number or pair of numbers that fixes the position of a point.}
\dictentry{cosine}{in a three-sided shape with a right angle, the ratio of the adjacent side to the hypotenuse.}
\dictentry{cosine rule}{a formula connecting three side lengths and one angle in a three-sided shape.}
\dictentry{counterexample}{an example that proves a general statement is not always true}
\dictentry{coupon}{The regular interest payment made on a bond.}
\dictentry{cube}{the number you get when a value is multiplied by itself twice.}
\dictentry{cuboid}{a 3D solid shape whose six faces are all rectangles.}
\dictentry{dataset}{a collection of numbers or information being studied}
\dictentry{decimal place}{the position of a digit after the decimal point, such as tenths or hundredths}
\dictentry{decimal places}{the digits written after the point in a number.}
\dictentry{decimal point}{the dot used to separate the whole part of a number from the fractional part.}
\dictentry{degree}{a unit for measuring angles; 360 of them make one full turn.}
\dictentry{depreciation}{the loss in value of something over time.}
\dictentry{dimension}{a measurement of a shape, such as its length or width.}
\dictentry{dimensions}{the measurements of a shape, such as length, width and height.}
\dictentry{direct proportion}{a relationship where two amounts change together while their ratio stays the same.}
\dictentry{directed path}{a line with an arrow showing direction from one place to another}
\dictentry{directly proportional}{when one amount is multiplied by a number, the other amount is multiplied by that same number.}
\dictentry{discriminant}{a value calculated from the numbers in an equation; if it is negative, the equation has no real solutions.}
\dictentry{displacement}{the straight-line distance and direction from a starting point to a final position.}
\dictentry{distributive property}{the rule that a(b+c)=ab+ac.}
\dictentry{divide}{to split a number into equal parts, or find how many times one number fits inside another.}
\dictentry{division}{the operation of splitting a number into equal parts.}
\dictentry{divisor}{the number that another number is divided by.}
\dictentry{duration}{the length of time until a bond or gilt is paid back.}
\dictentry{embedding}{a way of turning words into coordinates or vectors so meanings can be compared mathematically}
\dictentry{equal}{having the same size or number}
\dictentry{equation}{a mathematical statement with an equals sign showing that two sides are equal.}
\dictentry{equilateral}{having all sides the same length}
\dictentry{equilateral triangle}{a triangle with all three sides the same length.}
\dictentry{equivalent fraction}{a different fraction that means the same amount as another fraction.}
\dictentry{evaluate}{to work out the value of an expression}
\dictentry{even}{able to be divided exactly by 2.}
\dictentry{exponent}{the raised number showing how many times to multiply a number by itself.}
\dictentry{fibonacci}{a famous sequence where each number is the total of the two before it}
\dictentry{finite}{limited in size; not going on forever.}
\dictentry{fixed point}{a value that stays the same when an iterative process is applied to it.}
\dictentry{chunking}{a division method that breaks a larger number into easier smaller parts.}
\dictentry{cos}{the adjacent side divided by the hypotenuse for an angle in a right triangle.}
\dictentry{elimination}{removing one unknown from simultaneous equations by adding or subtracting them}
\dictentry{endpoint}{the point at which a route or line finishes}
\dictentry{even number}{a whole number that can be divided exactly by two.}
\dictentry{foil}{a way to multiply two brackets that each have two terms, multiplying First, Outer, Inner and Last.}
\dictentry{fraction wall}{a diagram of strips divided into equal parts, used to compare fractions.}
\dictentry{geometric}{describing a sequence where each term is multiplied by the same number to get the next.}
\dictentry{geometry}{the branch of mathematics that studies shapes, sizes, and space.}
\dictentry{gilt}{A bond issued by the UK government to borrow money.}
\dictentry{height}{the distance straight up from the base of a shape to its top, used to find area.}
\dictentry{hexagon}{a flat shape with six straight sides}
\dictentry{hundredths}{the value of one part when a whole is divided into one hundred equal parts.}
\dictentry{improper fraction}{a fraction with a top number larger than or equal to its bottom number}
\dictentry{interest}{Money paid for borrowing or earned on savings.}
\dictentry{interest rate}{the percentage of an amount that is paid or earned as interest.}
\dictentry{intersect}{to cross each other}
\dictentry{inverse sine}{The operation that finds which corner has a given sine value.}
\dictentry{isosceles triangle}{a three-sided shape with two sides equal in length}
\dictentry{iteration}{one complete repeat of a repeating process}
\dictentry{iterative}{using a repeated rule, where each result is used to find the next one.}
\dictentry{iterative process}{a method that repeats a rule, using the previous answer to produce the next}
\dictentry{lattice method}{a written multiplication method that uses a grid to organise the calculation.}
\dictentry{leg}{one of the two shorter sides of a right-angled triangle.}
\dictentry{like terms}{terms that have the same letter or letters, or no letter, so they can be added or subtracted together.}
\dictentry{line segment}{a straight part of a line that starts at one point and ends at another.}
\dictentry{magnitude}{the size of a number without its sign.}
\dictentry{major arc}{the longer of the two parts of the circumference between two points on a circle.}
\dictentry{major segment}{the larger of the two parts of a circle cut off by a chord.}
\dictentry{mapping}{matching or sending each point of one shape to a point of another.}
\dictentry{mass}{the amount of matter in an object, measured in grams or kilograms.}
\dictentry{maturity}{the time when a bond is paid back.}
\dictentry{minor arc}{the shorter of the two parts of the circumference between two points on a circle.}
\dictentry{minor segment}{the smaller of the two parts of a circle cut off by a chord.}
\dictentry{mixed number}{a number made of a whole amount and a fraction, such as 3 1/2}
\dictentry{negative}{below zero; a number or term with a minus sign is below zero.}
\dictentry{negative number}{a number less than zero, written with a minus sign.}
\dictentry{non-reflex angle}{the smaller angle between two lines, measuring less than 180 degrees}
\dictentry{number line}{a straight line with numbers marked at equal distances, used to show value and order.}
\dictentry{odd number}{a whole number that cannot be divided exactly by two.}
\dictentry{apex}{the point of a triangle opposite the base, often its top point.}
\dictentry{arithmetic sequence}{a sequence where the same number is added or subtracted each time}
\dictentry{compound}{to earn interest on both the original amount and on interest already earned.}
\dictentry{corresponding}{matching sides or points in similar shapes.}
\dictentry{enlarge}{to make a shape bigger by multiplying every length by the same number.}
\dictentry{fractional part}{the part of a mixed number that is less than one, written as a fraction.}
\dictentry{geometric sequence}{a sequence where each term is found by multiplying the previous term by the same number}
\dictentry{ladder}{an investment strategy that buys bonds with different maturity dates so payments arrive at regular times.}
\dictentry{odd}{a whole number that cannot be split exactly into two equal whole numbers.}
\dictentry{of}{in fraction work, it means multiply: taking a fraction of an amount uses multiplication}
\dictentry{operation}{Something done to a number, such as adding, subtracting, multiplying or dividing.}
\dictentry{opposite}{in a three-sided shape with a right angle, the side across from a chosen corner.}
\dictentry{order}{to arrange numbers from smallest to largest, or from largest to smallest.}
\dictentry{origin}{the fixed point (0,0) on a coordinate grid.}
\dictentry{parentheses}{the round brackets ( ) used in calculations to show which part to do first.}
\dictentry{per cent}{for every hundred; a number out of one hundred.}
\dictentry{percent}{for every hundred; shown with a \% sign.}
\dictentry{pi}{the special number about 3.14 used when working with circles}
\dictentry{pie chart}{a circular chart divided into slices to show parts of a whole.}
\dictentry{position-to-term rule}{a way to calculate any term in a sequence from its position, such as 3n+1.}
\dictentry{positive}{a number greater than zero.}
\dictentry{prime number}{a whole number greater than 1 whose only factor is 1 and itself.}
\dictentry{principal}{The original amount of money borrowed or lent, before interest or extra payments.}
\dictentry{proof}{a logical argument showing that a mathematical statement is always true.}
\dictentry{proper fraction}{a fraction with a top number smaller than its bottom number}
\dictentry{protractor}{a tool used to measure angles.}
\dictentry{pythagoras}{a rule for finding a missing side of a right-angled triangle.}
\dictentry{pythagoras' theorem}{in a right-angled triangle, the longest side times itself equals the other two sides times themselves, added together.}
\dictentry{quarter}{one of four equal parts of one.}
\dictentry{rate}{a comparison of two different kinds of amount, such as litres per minute.}
\dictentry{rate of change}{how quickly a measurement increases or decreases over time.}
\dictentry{ray}{a straight line that starts at one point and goes on forever in one direction}
\dictentry{rectangle}{a flat shape with four straight sides and four square corners.}
\dictentry{recurrence}{a rule that defines each term using one or more earlier terms}
\dictentry{recurring decimal}{a decimal in which a digit or group of digits repeats forever.}
\dictentry{regular}{having all sides equal and all angles equal}
\dictentry{remainder}{the amount left over when a number cannot be split into equal whole groups}
\dictentry{resultant}{the single vector that has the same overall effect as adding two or more vectors.}
\dictentry{return}{The amount of money gained or lost on an amount lent or invested.}
\dictentry{reverse}{The same vector with every number changed to the opposite sign, so it points back.}
\dictentry{right angle}{a square corner, exactly a quarter of a full turn.}
\dictentry{right triangle}{a triangle with one angle that is a quarter turn.}
\dictentry{right-angled triangle}{a three-sided shape with one corner of 90 degrees.}
\dictentry{rule}{a mathematical instruction that works out every number in a sequence}
\dictentry{scale}{to enlarge or shrink a shape so every length changes in the same way}
\dictentry{semicircular}{shaped like half of a round object.}
\dictentry{set}{a collection of numbers or objects grouped together.}
\dictentry{sign}{The plus or minus mark that tells whether a value is positive or negative.}
\dictentry{simple interest}{interest calculated only on the original amount, so previous interest earns nothing.}
\dictentry{simplest form}{a fraction written with the smallest possible numbers above and below the line.}
\dictentry{simultaneous equations}{two or more equations with the same unknown values that are solved together.}
\dictentry{sin}{the opposite side divided by the hypotenuse for an angle in a right triangle.}
\dictentry{sine}{in a right-angled triangle, the side across from a chosen corner divided by the longest side.}
\dictentry{sine rule}{a formula connecting the sides of a triangle with the sin of the angle opposite each side.}
\dictentry{solution}{a value or set of values that makes an equation true when used in place of the letter.}
\dictentry{speed}{how far something travels in one unit of time}
\dictentry{square}{to work out a number times itself; also a flat shape with four equal sides and four equal corners}
\dictentry{square number}{a result of multiplying a whole number by itself, such as 1, 4, 9, 16.}
\dictentry{subtended}{formed when the two ends of a chord are joined to a point on the circumference.}
\dictentry{subtraction}{the process of taking one amount away from another.}
\dictentry{symmetric}{balanced so that one side is a mirror image of the other.}
\dictentry{tan}{the opposite side divided by the adjacent side for an angle in a right triangle.}
\dictentry{ten-thousandths}{the value of one part when a whole is divided into ten thousand equal parts.}
\dictentry{tenths}{the value of one part when a whole is divided into ten equal parts.}
\dictentry{terminate}{for a decimal to end after a fixed number of digits.}
\dictentry{theorem}{a mathematical statement that can be proved to be true.}
\dictentry{thousandths}{the value of one part when a whole is divided into one thousand equal parts.}
\dictentry{times table}{a list of numbers made by repeatedly adding the same number, such as 3, 6, 9, 12.}
\dictentry{tip-to-tail}{A method for adding vectors by placing the point of one at the start of the next.}
\dictentry{total}{the whole amount found by adding all the parts together.}
\dictentry{translation}{moving every point by the same vector, without turning or resizing it}
\dictentry{triangle}{a flat shape with three straight sides}
\dictentry{triangular number}{a number that can be shown as dots arranged in a triangle}
\dictentry{trigonometry}{the mathematics of relationships between angles and side lengths in triangles.}
\dictentry{uniform}{the same at all points; not changing.}
\dictentry{unique}{being the only one; here, no other value is tied with the mode.}
\dictentry{unit price}{the cost for one single gram, litre, item, or other measure}
\dictentry{upper bound}{a number that is greater than or equal to every possible value in a set.}
\dictentry{value}{the number that a letter or symbol stands for}
\dictentry{vertical}{going straight up and down}
\dictentry{vertically opposite}{describing two angles, the same size, formed across from each other when two lines cross}
\dictentry{whole number}{a number that is not negative and has no part after a decimal point, such as 0, 1, 2 or 3.}
\dictentry{whole numbers}{the numbers zero, one, two, three, and so on, with no fraction or decimal part.}
\dictentry{yield}{The money earned on an amount lent or invested, usually shown as a percentage.}
\dictentry{yield curve}{a line showing how yields change for different times to maturity.}

\end{multicols}

\end{document}
